
\subsubsection{2.1 Neural Network Calibration}

Guo et al. (2017) demonstrated that modern neural networks are often miscalibrated and introduced temperature scaling as a simple post-hoc correction that reduces ECE. This work has been highly influential (7,452 citations as of the search date) and established temperature scaling as a standard baseline for calibration methods.

Recent work has extended calibration methods to large language models. DACA \citep{Luo2025DACA} uses disagreement between pre-trained and post-trained language models to detect samples requiring calibration, reporting 15\% ECE improvement on Llama-3. CCPS \citep{Khanmohammadi2025CCPS} uses consistency under perturbation as a calibration signal, reporting 55\% ECE reduction on MMLU. These methods focus on ECE as the primary metric and do not directly measure discriminative quality.

\subsubsection{2.2 RLHF and Overconfidence}

Tian et al. (2023) documented that RLHF-tuned models exhibit systematic overconfidence and found that verbalized confidence often outperforms token-level probabilities for such models. This suggests that RLHF training affects how models express uncertainty at the probability level.

The Bradley-Terry preference model underlying RLHF training rewards responses preferred by human annotators. To the extent that annotators prefer confident-sounding responses, this creates selection pressure toward higher confidence regardless of correctness.

\subsubsection{2.3 Brier Score Decomposition}

Murphy (1973) introduced a decomposition of the Brier score into reliability, resolution (refinement), and uncertainty components. This framework separates calibration error (reliability) from discriminative ability (refinement), enabling characterization of whether miscalibration is primarily a scale effect or a shape effect.
